\chapter{The Conservation of Resolution: Gravity as an Informational Knot}

In the previous chapter we derived a single integer --- $n \approx 184$
bits --- that fixes the resolution of space and time for our universe.
That derivation treated the universe as an empty system, a pure
information budget with no internal structure. Now we ask what happens
when matter enters the picture. 

We start with only one principle, already implicit in
everything that came before: the bit budget is fixed. 

\section{The Knot in the Rope}

Think of the universe's 184-bit budget as a length of rope. When the
rope is smooth and uniform, every segment is independently addressable.
An observer running their fingers along it encounters distinct positions,
one after another.

Now tie a knot. The material of the rope is conserved --- nothing has
been added or removed. But a section of the rope has been drawn inward
and bundled. If you stretch the rope out, its total navigable length has
shrunk. The bits caught in the knot are no longer available as
independently addressable positions along the background.

This is exactly what happens when matter forms. At any moment, the $n$
bits of the universe are partitioned between two roles:

\begin{itemize}
    \item \textbf{Spacetime fabric}: bits not consumed by any composite
    structure, remaining independently addressable as background geometry.
    \item \textbf{Matter structures}: bits bundled into composite objects
    at successive levels --- hadrons, atoms, compounds --- each consuming
    a fixed number of fabric tokens from the available pool.
\end{itemize}

Because no new bits can be injected from outside, this is a strict
zero-sum partition. If $m(t)$ bits have been consumed by matter
structures at moment $t$, then

\[
\rho_{\mathrm{fabric}}(t) = n - m(t).
\]

Every bit locked into a particle is a bit subtracted from the spatial
fabric available to an internal observer.

\section{The Relational Scale Factor}

An observer inside the universe cannot step outside and measure it with
an external ruler. They can only count the structures available to them:
the free fabric tokens and the composite matter entities. The resolution
of observable space --- the number of independently distinguishable
positions --- is the total of both.

When one composite structure of width $w$ bits forms, it consumes $w$
fabric tokens and creates $1$ matter entity. The new resolution is:

\[
(n - w) + 1 = n - (w - 1).
\]

Resolution decreases by $w - 1$ for every structure that forms. If
$m(t)$ bits have been consumed and $k(t)$ composite entities exist, the
total resolution is:

\[
\mathrm{Resolution}(t) = \bigl(n - m(t)\bigr) + k(t).
\]

Normalised to the maximum possible resolution of an empty universe, this
defines the relational scale factor:

\begin{equation}
R(t) = \frac{\bigl(n - m(t)\bigr) + k(t)}{n}.
\label{eq:scalefactor}
\end{equation}

This is a pure counting statement which does not need metric tensor, field equation 
or force law. Just: what fraction of the bit budget remains independently
addressable to an internal observer?

\section{Two Famous Solutions from One Equation}

The two most celebrated exact solutions of General Relativity emerge
directly as the extreme values of equation~\eqref{eq:scalefactor}.

\textbf{The De Sitter limit.} Set $m = 0$ and $k = 0$: no matter has
formed, all bits are free fabric. Then $R = 1$. The spatial grid retains
its maximum addressable resolution. An internal observer perceives pure
exponential expansion without bound --- exactly the De Sitter vacuum
solution of GR with a positive cosmological constant and zero matter
density.

\textbf{The Schwarzschild limit.} Set $m = n$ and $k = 1$: every bit in
the universe has been consumed into a single composite entity. Then:

\[
R = \frac{0 + 1}{n} = \frac{1}{n}.
\]

The observable resolution of space has collapsed to a single isolated
coordinate. Space has effectively vanished. This is the informational
counterpart of the black hole singularity established in Chapter~2 ---
and it arrives here not as an assumption but as the natural far end of
the same counting equation that gives De Sitter at the other extreme.

These are not two separate mathematical discoveries about different
aspects of gravity. They are the two endpoints of a single
information-conservation spectrum. Every state of the observed universe
--- expanding, matter-filled, approaching equilibrium --- is a point
somewhere between them.

\section{Re-Reading the Friedmann Equation}

The Friedmann equation governs the expansion of a homogeneous, isotropic
universe in General Relativity:

\[
\left(\frac{\dot{a}}{a}\right)^2 = \frac{8\pi G}{3}\rho_{\mathrm{matter}}
+ \frac{\Lambda}{3}.
\]

The present framework maps onto this equation under identifications that
follow directly from the counting argument:

\[
\rho_{\mathrm{matter}} \;\longleftrightarrow\; \frac{k(t)}{n},
\qquad
\Lambda \;\longleftrightarrow\; \frac{\rho_{\mathrm{fabric}}(t)}{n},
\qquad
a \;\longleftrightarrow\; R(t).
\]

These are not free parameters fitted to observation. They are
definitions. Matter density is the fraction of the budget currently
locked into bound entities. The cosmological constant is the fraction
that remains free. The scale factor is the addressable resolution
fraction.

Read this way, the Friedmann equation is not a dynamical law imposed on
spacetime. It is a counting identity, a statement about how the bit
budget is partitioned at each moment in the entropy evolution of the
universe.

\section{The Cosmological Constant Problem Dissolved}

Quantum field theory has predicted a vacuum energy density
roughly $10^{120}$ times larger than the observed value of $\Lambda$.
This is often called the worst prediction in the history of physics.

The framework dissolves the problem by removing its premise. The
cosmological constant is not a property of the vacuum in this picture.
It is not an energy density of empty space at all. It is the fraction of
the bit budget that happens to be unbound at this cosmic epoch --- a
dimensionless ratio between zero and one, bounded above by the finite
integer $n$.

There is no infinite vacuum energy to cancel. There is only a fabric
density that grows when matter decays and shrinks when matter forms. The
$10^{120}$ discrepancy is a consequence of mistaking a counting ratio
for an energy density. The quantities are not the same thing.

\section{Gravity as Geometry, Not Force}

In General Relativity, matter curves spacetime through the stress-energy
tensor --- matter tells spacetime how to curve, spacetime tells matter
how to move. The mechanism is geometrically precise but physically
unexplained. Why does matter curve space at all?

In this framework the answer is direct. A composite structure consumes
fabric tokens from the surrounding region. The local density of
independently addressable positions drops. These structures are
building blocks of ever larger structures, ultimately the observer.
Probability gradient points towards higher knot density. However,
the higher knot density means fewer and fewer resolvable positions.
This depletion of local resolution is what they
experience as gravitational attraction. Space curves near a mass because
the mass has consumed the very bits needed to render the background
between them.

Gravity is not a force thrown across empty space. It is the geometric
consequence of an informational redistribution --- the same knot-in-a-rope
phenomenon, read from the inside.

\section{From Linear Counting to Spherical Geometry}

The counting equation \eqref{eq:scalefactor} is intrinsically
one-dimensional: it counts bits along a linear budget. The observed
universe is three-dimensional and spherically symmetric. The bridge
between the two is the same $4\pi$ factor that appeared in Chapter~2
when recovering the Bekenstein--Hawking entropy:

\[
a_{\mathrm{obs}} \propto 4\pi \cdot R(t)^2.
\]

Projecting the linear bit-depletion curve through a spherical surface
converts it into the quadratic density profile found inside the event
horizon of a Schwarzschild black hole. The geometry follows the data.

The $4\pi$ factor is the unique
conversion between a one-dimensional radial information count and a
three-dimensional spherical observable, the same conversion that
appeared in Chapter~2 and will reappear whenever the framework
transitions between informational and geometric descriptions.

It should be noted that the $4\pi$ conversion factor was not
strictly derived from $n$. We might ask: why don't we live on an
open-edged flat plane? The answer could be that a 3D sphere is the
most straightforward geometry that prevents us from walking off the
edge and falling into... where?

As the answer is terrifying, lets fix it by introducing $4\pi$.


\section{What This Chapter Establishes}

Two results follow from the information-conservation argument.

First, De Sitter space and the Schwarzschild singularity are the two
extreme values of a single counting equation. De Sitter corresponds to
all bits free; Schwarzschild corresponds to all bits bound into one
entity. General Relativity's two boundary solutions are not independent
discoveries. They are the endpoints of one information-theoretic
spectrum.

Second, the Friedmann equation maps onto the framework under explicit
variable identifications that follow from the counting argument alone.
The cosmological constant is the free fabric fraction. Matter density is
the bound entity fraction. The scale factor is the addressable resolution
fraction. No constants are fitted. No forces are introduced.

What the framework has not yet explained is why individual particles,
moving through this information grid, behave the way they do at small
scales --- why their trajectories are probabilistic, why they interfere,
and why measurement produces a definite outcome. That is the question the
next chapter answers. The same spectral complexity functional that
resolved the singularity and selected Keplerian orbits will now be shown
to produce quantum mechanics.
